\IfFileExists{stacks-project.cls}{%
\documentclass{stacks-project}
}{%
\documentclass{amsart}
}

% For dealing with references we use the comment environment
\usepackage{verbatim}
\newenvironment{reference}{\comment}{\endcomment}
%\newenvironment{reference}{}{}
\newenvironment{slogan}{\comment}{\endcomment}
\newenvironment{history}{\comment}{\endcomment}

% For commutative diagrams we use Xy-pic
\usepackage[all]{xy}

% We use 2cell for 2-commutative diagrams.
\xyoption{2cell}
\UseAllTwocells

% We use multicol for the list of chapters between chapters
\usepackage{multicol}

% This is generally recommended for better output
\usepackage{lmodern}
\usepackage[T1]{fontenc}

% For cross-file-references
\usepackage{xr-hyper}

% Package for hypertext links:
\usepackage{hyperref}

% For any local file, say "hello.tex" you want to link to please
% use \externaldocument[hello-]{hello}
\externaldocument[introduction-]{introduction}
\externaldocument[conventions-]{conventions}
\externaldocument[sets-]{sets}
\externaldocument[categories-]{categories}
\externaldocument[topology-]{topology}
\externaldocument[sheaves-]{sheaves}
\externaldocument[sites-]{sites}
\externaldocument[stacks-]{stacks}
\externaldocument[fields-]{fields}
\externaldocument[algebra-]{algebra}
\externaldocument[brauer-]{brauer}
\externaldocument[homology-]{homology}
\externaldocument[derived-]{derived}
\externaldocument[simplicial-]{simplicial}
\externaldocument[more-algebra-]{more-algebra}
\externaldocument[smoothing-]{smoothing}
\externaldocument[modules-]{modules}
\externaldocument[sites-modules-]{sites-modules}
\externaldocument[injectives-]{injectives}
\externaldocument[cohomology-]{cohomology}
\externaldocument[sites-cohomology-]{sites-cohomology}
\externaldocument[dga-]{dga}
\externaldocument[dpa-]{dpa}
\externaldocument[sdga-]{sdga}
\externaldocument[hypercovering-]{hypercovering}
\externaldocument[schemes-]{schemes}
\externaldocument[constructions-]{constructions}
\externaldocument[properties-]{properties}
\externaldocument[morphisms-]{morphisms}
\externaldocument[coherent-]{coherent}
\externaldocument[divisors-]{divisors}
\externaldocument[limits-]{limits}
\externaldocument[varieties-]{varieties}
\externaldocument[topologies-]{topologies}
\externaldocument[descent-]{descent}
\externaldocument[perfect-]{perfect}
\externaldocument[more-morphisms-]{more-morphisms}
\externaldocument[flat-]{flat}
\externaldocument[groupoids-]{groupoids}
\externaldocument[more-groupoids-]{more-groupoids}
\externaldocument[etale-]{etale}
\externaldocument[chow-]{chow}
\externaldocument[intersection-]{intersection}
\externaldocument[pic-]{pic}
\externaldocument[weil-]{weil}
\externaldocument[adequate-]{adequate}
\externaldocument[dualizing-]{dualizing}
\externaldocument[duality-]{duality}
\externaldocument[discriminant-]{discriminant}
\externaldocument[derham-]{derham}
\externaldocument[local-cohomology-]{local-cohomology}
\externaldocument[algebraization-]{algebraization}
\externaldocument[curves-]{curves}
\externaldocument[resolve-]{resolve}
\externaldocument[models-]{models}
\externaldocument[functors-]{functors}
\externaldocument[equiv-]{equiv}
\externaldocument[pione-]{pione}
\externaldocument[etale-cohomology-]{etale-cohomology}
\externaldocument[proetale-]{proetale}
\externaldocument[relative-cycles-]{relative-cycles}
\externaldocument[more-etale-]{more-etale}
\externaldocument[trace-]{trace}
\externaldocument[crystalline-]{crystalline}
\externaldocument[spaces-]{spaces}
\externaldocument[spaces-properties-]{spaces-properties}
\externaldocument[spaces-morphisms-]{spaces-morphisms}
\externaldocument[decent-spaces-]{decent-spaces}
\externaldocument[spaces-cohomology-]{spaces-cohomology}
\externaldocument[spaces-limits-]{spaces-limits}
\externaldocument[spaces-divisors-]{spaces-divisors}
\externaldocument[spaces-over-fields-]{spaces-over-fields}
\externaldocument[spaces-topologies-]{spaces-topologies}
\externaldocument[spaces-descent-]{spaces-descent}
\externaldocument[spaces-perfect-]{spaces-perfect}
\externaldocument[spaces-more-morphisms-]{spaces-more-morphisms}
\externaldocument[spaces-flat-]{spaces-flat}
\externaldocument[spaces-groupoids-]{spaces-groupoids}
\externaldocument[spaces-more-groupoids-]{spaces-more-groupoids}
\externaldocument[bootstrap-]{bootstrap}
\externaldocument[spaces-pushouts-]{spaces-pushouts}
\externaldocument[spaces-chow-]{spaces-chow}
\externaldocument[groupoids-quotients-]{groupoids-quotients}
\externaldocument[spaces-more-cohomology-]{spaces-more-cohomology}
\externaldocument[spaces-simplicial-]{spaces-simplicial}
\externaldocument[spaces-duality-]{spaces-duality}
\externaldocument[formal-spaces-]{formal-spaces}
\externaldocument[restricted-]{restricted}
\externaldocument[spaces-resolve-]{spaces-resolve}
\externaldocument[formal-defos-]{formal-defos}
\externaldocument[defos-]{defos}
\externaldocument[cotangent-]{cotangent}
\externaldocument[examples-defos-]{examples-defos}
\externaldocument[algebraic-]{algebraic}
\externaldocument[examples-stacks-]{examples-stacks}
\externaldocument[stacks-sheaves-]{stacks-sheaves}
\externaldocument[criteria-]{criteria}
\externaldocument[artin-]{artin}
\externaldocument[quot-]{quot}
\externaldocument[stacks-properties-]{stacks-properties}
\externaldocument[stacks-morphisms-]{stacks-morphisms}
\externaldocument[stacks-limits-]{stacks-limits}
\externaldocument[stacks-cohomology-]{stacks-cohomology}
\externaldocument[stacks-perfect-]{stacks-perfect}
\externaldocument[stacks-introduction-]{stacks-introduction}
\externaldocument[stacks-more-morphisms-]{stacks-more-morphisms}
\externaldocument[stacks-geometry-]{stacks-geometry}
\externaldocument[moduli-]{moduli}
\externaldocument[moduli-curves-]{moduli-curves}
\externaldocument[examples-]{examples}
\externaldocument[exercises-]{exercises}
\externaldocument[guide-]{guide}
\externaldocument[desirables-]{desirables}
\externaldocument[coding-]{coding}
\externaldocument[obsolete-]{obsolete}
\externaldocument[fdl-]{fdl}
\externaldocument[index-]{index}

% Theorem environments.
%
\theoremstyle{plain}
\newtheorem{theorem}[subsection]{Theorem}
\newtheorem{proposition}[subsection]{Proposition}
\newtheorem{lemma}[subsection]{Lemma}

\theoremstyle{definition}
\newtheorem{definition}[subsection]{Definition}
\newtheorem{example}[subsection]{Example}
\newtheorem{exercise}[subsection]{Exercise}
\newtheorem{situation}[subsection]{Situation}

\theoremstyle{remark}
\newtheorem{remark}[subsection]{Remark}
\newtheorem{remarks}[subsection]{Remarks}

\numberwithin{equation}{subsection}

% Macros
%
\def\lim{\mathop{\mathrm{lim}}\nolimits}
\def\colim{\mathop{\mathrm{colim}}\nolimits}
\def\Spec{\mathop{\mathrm{Spec}}}
\def\Hom{\mathop{\mathrm{Hom}}\nolimits}
\def\Ext{\mathop{\mathrm{Ext}}\nolimits}
\def\SheafHom{\mathop{\mathcal{H}\!\mathit{om}}\nolimits}
\def\SheafExt{\mathop{\mathcal{E}\!\mathit{xt}}\nolimits}
\def\Sch{\mathit{Sch}}
\def\Mor{\mathop{\mathrm{Mor}}\nolimits}
\def\Ob{\mathop{\mathrm{Ob}}\nolimits}
\def\Sh{\mathop{\mathit{Sh}}\nolimits}
\def\NL{\mathop{N\!L}\nolimits}
\def\CH{\mathop{\mathrm{CH}}\nolimits}
\def\proetale{{pro\text{-}\acute{e}tale}}
\def\etale{{\acute{e}tale}}
\def\QCoh{\mathit{QCoh}}
\def\Ker{\mathop{\mathrm{Ker}}}
\def\Im{\mathop{\mathrm{Im}}}
\def\Coker{\mathop{\mathrm{Coker}}}
\def\Coim{\mathop{\mathrm{Coim}}}

% Boxtimes
%
\DeclareMathSymbol{\boxtimes}{\mathbin}{AMSa}{"02}

%
% Macros for moduli stacks/spaces
%
\def\QCohstack{\mathcal{QC}\!\mathit{oh}}
\def\Cohstack{\mathcal{C}\!\mathit{oh}}
\def\Spacesstack{\mathcal{S}\!\mathit{paces}}
\def\Quotfunctor{\mathrm{Quot}}
\def\Hilbfunctor{\mathrm{Hilb}}
\def\Curvesstack{\mathcal{C}\!\mathit{urves}}
\def\Polarizedstack{\mathcal{P}\!\mathit{olarized}}
\def\Complexesstack{\mathcal{C}\!\mathit{omplexes}}
% \Pic is the operator that assigns to X its picard group, usage \Pic(X)
% \Picardstack_{X/B} denotes the Picard stack of X over B
% \Picardfunctor_{X/B} denotes the Picard functor of X over B
\def\Pic{\mathop{\mathrm{Pic}}\nolimits}
\def\Picardstack{\mathcal{P}\!\mathit{ic}}
\def\Picardfunctor{\mathrm{Pic}}
\def\Deformationcategory{\mathcal{D}\!\mathit{ef}}
\usepackage{cleveref}
\crefname{originallemma}{original lemma}{original lemmas}
\Crefname{originallemma}{Original lemma}{Original lemmas}
\makeatletter
\newlabel{cohomology-lemma-mayer-vietoris}{{01EB}{source}{}{sourcepointer-1}{}}
\newlabel{cohomology-lemma-mayer-vietoris@cref}{{[originallemma][0][]01EB}{[1][0][]source}{}{}{}}
\makeatother
\crefname{originallemma}{original lemma}{original lemmas}
\Crefname{originallemma}{Original lemma}{Original lemmas}
\makeatletter
\newlabel{lemma-cokernel-adequate}{{06V2}{source}{}{sourcepointer-2}{}}
\newlabel{lemma-cokernel-adequate@cref}{{[originallemma][0][]06V2}{[1][0][]source}{}{}{}}
\makeatother
\crefname{originallemma}{original lemma}{original lemmas}
\Crefname{originallemma}{Original lemma}{Original lemmas}
\makeatletter
\newlabel{lemma-extension-adequate}{{06V8}{source}{}{sourcepointer-3}{}}
\newlabel{lemma-extension-adequate@cref}{{[originallemma][0][]06V8}{[1][0][]source}{}{}{}}
\makeatother
\crefname{originallemma}{original lemma}{original lemmas}
\Crefname{originallemma}{Original lemma}{Original lemmas}
\makeatletter
\newlabel{lemma-kernel-adequate}{{06V3}{source}{}{sourcepointer-4}{}}
\newlabel{lemma-kernel-adequate@cref}{{[originallemma][0][]06V3}{[1][0][]source}{}{}{}}
\makeatother
\crefname{originallemma}{original lemma}{original lemmas}
\Crefname{originallemma}{Original lemma}{Original lemmas}
\makeatletter
\newlabel{lemma-pushforward-adequate}{{06VA}{source}{}{sourcepointer-5}{}}
\newlabel{lemma-pushforward-adequate@cref}{{[originallemma][0][]06VA}{[1][0][]source}{}{}{}}
\makeatother
\crefname{originallemma}{original lemma}{original lemmas}
\Crefname{originallemma}{Original lemma}{Original lemmas}
\makeatletter
\newlabel{lemma-same-cohomology-adequate}{{06VM}{source}{}{sourcepointer-6}{}}
\newlabel{lemma-same-cohomology-adequate@cref}{{[originallemma][0][]06VM}{[1][0][]source}{}{}{}}
\makeatother
\crefname{originallemma}{original lemma}{original lemmas}
\Crefname{originallemma}{Original lemma}{Original lemmas}
\makeatletter
\newlabel{lemma-sheafification-adequate}{{06VP}{source}{}{sourcepointer-7}{}}
\newlabel{lemma-sheafification-adequate@cref}{{[originallemma][0][]06VP}{[1][0][]source}{}{}{}}
\makeatother
\crefname{originallemma}{original lemma}{original lemmas}
\Crefname{originallemma}{Original lemma}{Original lemmas}
\makeatletter
\newlabel{schemes-lemma-characterize-separated}{{01KP}{source}{}{sourcepointer-8}{}}
\newlabel{schemes-lemma-characterize-separated@cref}{{[originallemma][0][]01KP}{[1][0][]source}{}{}{}}
\makeatother
\crefname{originallemma}{original lemma}{original lemmas}
\Crefname{originallemma}{Original lemma}{Original lemmas}
\makeatletter
\newlabel{sites-cohomology-lemma-cohomology-of-open}{{03F3}{source}{}{sourcepointer-9}{}}
\newlabel{sites-cohomology-lemma-cohomology-of-open@cref}{{[originallemma][0][]03F3}{[1][0][]source}{}{}{}}
\makeatother
\crefname{originallemma}{original lemma}{original lemmas}
\Crefname{originallemma}{Original lemma}{Original lemmas}
\makeatletter
\newlabel{topologies-lemma-morphism-big-fppf}{{021W}{source}{}{sourcepointer-10}{}}
\newlabel{topologies-lemma-morphism-big-fppf@cref}{{[originallemma][0][]021W}{[1][0][]source}{}{}{}}
\makeatother
\title{Pushforwards and higher direct images of adequate modules are adequate}
\author{The Stacks Project Authors}
\date{stacks-project:adequate.tex:local-snapshot-20261012; GNU FDL 1.2 or later}
\begin{document}
\maketitle
\noindent\textit{Editorial note.} All delivered mathematical text is the author's own, copied from the pinned Stacks Project chapter file adequate.tex; no statement or proof step was rewritten, reordered or removed.\par\medskip
\noindent\textit{Editorial note.} The proof's closing sentence states that the last case 'proceeds in exactly the same manner as in the paragraph above and is omitted'; this is the author's own named-omission of a repetition of the immediately preceding displayed induction step, and the sentence is delivered verbatim rather than summarised.\par\medskip
\noindent\textit{Editorial note.} The counted spans carry navigation stubs for the references leaving them; no bibliography citation occurs in the spans.\par\medskip
\noindent\textit{Editorial note.} The preamble span points at the chapter's own preamble.tex shipped with the frozen source, used verbatim as staged.\par\medskip
\noindent\textit{Allowed standard prerequisites (editorial).}\begin{itemize}
\item Computation of the cohomology of an object of a big site, and of sheaves on a site as cohomology of the small site. Topologies, Lemma topologies-lemma-morphism-big-fppf and Cohomology on Sites, Lemma sites-cohomology-lemma-cohomology-of-open, invoked at proof lines 2138-2146. Give the first and third equalities in the computation of R\textasciicircum{}i f\_* F.
\item Identification of the cohomology of an adequate module computed in the big and in the Zariski topology. Same chapter, Lemma lemma-same-cohomology-adequate, invoked at proof lines 2146 and 2156-2160. Gives the last equality and identifies F\_B as an ordinary quasi-coherent sheaf.
\item Adequacy of a functor is detected after sheafification. Same chapter, Lemma lemma-sheafification-adequate, invoked at proof lines 2147-2149. Reduces the claim to the displayed claim about the functor B -> H\textasciicircum{}q(T\_B, F\_B).
\item Adequacy of the zeroth cohomology functor for affine T. Same chapter, Lemma lemma-pushforward-adequate, invoked at proof lines 2171-2172. Settles Case I.
\end{itemize}\medskip
\ifdefined\hypertarget\hypertarget{sourcepointer-1}{}\fi
\noindent\textit{Source reference.} Original lemma 01EB: sources/stacks-chow/tags.txt, source lines 1828--1828. Pointer outside the counted span, invoked once for context; the load-bearing argument of the counted statement is the author's proof reproduced in full in the proof span.\par\medskip
\ifdefined\hypertarget\hypertarget{sourcepointer-2}{}\fi
\noindent\textit{Source reference.} Original lemma 06V2: sources/stacks-adequate/pinned-original-source.tex, source lines 684--684. Pointer outside the counted span, invoked once for context; the load-bearing argument of the counted statement is the author's proof reproduced in full in the proof span.\par\medskip
\ifdefined\hypertarget\hypertarget{sourcepointer-3}{}\fi
\noindent\textit{Source reference.} Original lemma 06V8: sources/stacks-adequate/pinned-original-source.tex, source lines 910--910. Pointer outside the counted span, invoked once for context; the load-bearing argument of the counted statement is the author's proof reproduced in full in the proof span.\par\medskip
\ifdefined\hypertarget\hypertarget{sourcepointer-4}{}\fi
\noindent\textit{Source reference.} Original lemma 06V3: sources/stacks-adequate/pinned-original-source.tex, source lines 709--709. Pointer outside the counted span, invoked once for context; the load-bearing argument of the counted statement is the author's proof reproduced in full in the proof span.\par\medskip
\ifdefined\hypertarget\hypertarget{sourcepointer-5}{}\fi
\noindent\textit{Source reference.} Original lemma 06VA: sources/stacks-adequate/pinned-original-source.tex, source lines 968--968. Pointer outside the counted span, invoked once for context; the load-bearing argument of the counted statement is the author's proof reproduced in full in the proof span.\par\medskip
\ifdefined\hypertarget\hypertarget{sourcepointer-6}{}\fi
\noindent\textit{Source reference.} Original lemma 06VM: sources/stacks-adequate/pinned-original-source.tex, source lines 1896--1896. Pointer outside the counted span, invoked once for context; the load-bearing argument of the counted statement is the author's proof reproduced in full in the proof span.\par\medskip
\ifdefined\hypertarget\hypertarget{sourcepointer-7}{}\fi
\noindent\textit{Source reference.} Original lemma 06VP: sources/stacks-adequate/pinned-original-source.tex, source lines 1985--1985. Pointer outside the counted span, invoked once for context; the load-bearing argument of the counted statement is the author's proof reproduced in full in the proof span.\par\medskip
\ifdefined\hypertarget\hypertarget{sourcepointer-8}{}\fi
\noindent\textit{Source reference.} Original lemma 01KP: sources/stacks-chow/tags.txt, source lines 2058--2058. Pointer outside the counted span, invoked once for context; the load-bearing argument of the counted statement is the author's proof reproduced in full in the proof span.\par\medskip
\ifdefined\hypertarget\hypertarget{sourcepointer-9}{}\fi
\noindent\textit{Source reference.} Original lemma 03F3: sources/stacks-chow/tags.txt, source lines 4478--4478. Pointer outside the counted span, invoked once for context; the load-bearing argument of the counted statement is the author's proof reproduced in full in the proof span.\par\medskip
\ifdefined\hypertarget\hypertarget{sourcepointer-10}{}\fi
\noindent\textit{Source reference.} Original lemma 021W: sources/stacks-chow/tags.txt, source lines 2685--2685. Pointer outside the counted span, invoked once for context; the load-bearing argument of the counted statement is the author's proof reproduced in full in the proof span.\par\medskip
\begin{lemma}
\label{lemma-direct-image-adequate}
Let $f : T \to S$ be a quasi-compact and quasi-separated morphism
of schemes. For any adequate $\mathcal{O}_T$-module on
$(\Sch/T)_\tau$ the pushforward
$f_*\mathcal{F}$ and the higher direct images $R^if_*\mathcal{F}$
are adequate $\mathcal{O}_S$-modules on $(\Sch/S)_\tau$.
\end{lemma}

\begin{proof}
First we explain how to compute the higher direct images.
Choose an injective resolution $\mathcal{F} \to \mathcal{I}^\bullet$.
Then $R^if_*\mathcal{F}$ is the $i$th cohomology sheaf of the
complex $f_*\mathcal{I}^\bullet$.
Hence $R^if_*\mathcal{F}$ is the sheaf associated to the presheaf
which associates to an object $U/S$ of $(\Sch/S)_\tau$
the module
\begin{align*}
\frac{\Ker(f_*\mathcal{I}^i(U) \to f_*\mathcal{I}^{i + 1}(U))}
{\Im(f_*\mathcal{I}^{i - 1}(U) \to f_*\mathcal{I}^i(U))}
& =
\frac{\Ker(\mathcal{I}^i(U \times_S T) \to
\mathcal{I}^{i + 1}(U \times_S T))}
{\Im(\mathcal{I}^{i - 1}(U \times_S T) \to \mathcal{I}^i(U \times_S T))}
\\
& =
H^i(U \times_S T, \mathcal{F}) \\
& = H^i((\Sch/U \times_S T)_\tau,
\mathcal{F}|_{(\Sch/U \times_S T)_\tau}) \\
& = H^i(U \times_S T, \mathcal{F}|_{(U \times_S T)_{Zar}})
\end{align*}
The first equality by
Topologies, Lemma \ref{topologies-lemma-morphism-big-fppf}
(and its analogues for other topologies),
the second equality by definition of cohomology of $\mathcal{F}$
over an object of $(\Sch/T)_\tau$,
the third equality by
Cohomology on Sites, Lemma \ref{sites-cohomology-lemma-cohomology-of-open},
and the last equality by
Lemma \ref{lemma-same-cohomology-adequate}.
Thus by
Lemma \ref{lemma-sheafification-adequate}
it suffices to prove the claim stated in the following paragraph.

\medskip\noindent
Let $A$ be a ring. Let $T$ be a scheme quasi-compact and quasi-separated
over $A$. Let $\mathcal{F}$ be an adequate $\mathcal{O}_T$-module on
$(\Sch/T)_\tau$. For an $A$-algebra $B$ set
$T_B = T \times_{\Spec(A)} \Spec(B)$ and denote
$\mathcal{F}_B = \mathcal{F}|_{(T_B)_{Zar}}$ the restriction of
$\mathcal{F}$ to the small Zariski site of $T_B$.
(Recall that this is a ``usual'' quasi-coherent sheaf on the scheme
$T_B$, see
Lemma \ref{lemma-same-cohomology-adequate}.)
Claim: The functor
$$
B \longmapsto H^q(T_B, \mathcal{F}_B)
$$
is adequate. We will prove the lemma by the usual
procedure of cutting $T$ into pieces.

\medskip\noindent
Case I: $T$ is affine. In this case the schemes $T_B$ are all affine
and $H^q(T_B, \mathcal{F}_B) = 0$ for all $q \geq 1$.
The functor $B \mapsto H^0(T_B, \mathcal{F}_B)$ is adequate by
Lemma \ref{lemma-pushforward-adequate}.

\medskip\noindent
Case II: $T$ is separated. Let $n$ be the minimal number of affines needed
to cover $T$. We argue by induction on $n$. The base case is Case I.
Choose an affine open covering $T = V_1 \cup \ldots \cup V_n$.
Set $V = V_1 \cup \ldots \cup V_{n - 1}$ and $U = V_n$. Observe that
$$
U \cap V = (V_1 \cap V_n) \cup \ldots \cup (V_{n - 1} \cap V_n)
$$
is also a union of $n - 1$ affine opens as $T$ is separated, see
Schemes, Lemma \ref{schemes-lemma-characterize-separated}.
Note that for each $B$ the base changes $U_B$, $V_B$ and
$(U \cap V)_B = U_B \cap V_B$ behave in the same way. Hence we see that
for each $B$ we have a long exact sequence
$$
0 \to
H^0(T_B, \mathcal{F}_B) \to
H^0(U_B, \mathcal{F}_B) \oplus H^0(V_B, \mathcal{F}_B) \to
H^0((U \cap V)_B, \mathcal{F}_B) \to
H^1(T_B, \mathcal{F}_B) \to \ldots
$$
functorial in $B$, see
Cohomology, Lemma \ref{cohomology-lemma-mayer-vietoris}.
By induction hypothesis the functors
$B \mapsto H^q(U_B, \mathcal{F}_B)$,
$B \mapsto H^q(V_B, \mathcal{F}_B)$, and
$B \mapsto H^q((U \cap V)_B, \mathcal{F}_B)$
are adequate. Using
Lemmas \ref{lemma-kernel-adequate} and
\ref{lemma-cokernel-adequate}
we see that our functor $B \mapsto H^q(T_B, \mathcal{F}_B)$ sits in the
middle of a short exact sequence whose outer terms are adequate.
Thus the claim follows from
Lemma \ref{lemma-extension-adequate}.

\medskip\noindent
Case III: General quasi-compact and quasi-separated case.
The proof is again by induction on the number $n$ of affines needed to
cover $T$. The base case $n = 1$ is Case I.
Choose an affine open covering $T = V_1 \cup \ldots \cup V_n$.
Set $V = V_1 \cup \ldots \cup V_{n - 1}$ and $U = V_n$. Note that
since $T$ is quasi-separated $U \cap V$ is a quasi-compact open of an
affine scheme, hence Case II applies to it. The rest of the argument
proceeds in exactly the same manner as in the paragraph above and is
omitted.
\end{proof}

\end{document}
